\documentclass[sigconf,nonacm]{acmart}

% Citations need BibTeX: \texttt{references.bib} must sit beside \texttt{main.tex}.
% With \texttt{-output-directory=out}, run \texttt{bibtex out/main} after the first
% \texttt{pdflatex}, then \texttt{pdflatex} twice more (or \texttt{make build} / \texttt{latexmk}).

% ── Packages ──────────────────────────────────────────────────────────────────
\usepackage{listings}
\usepackage{xcolor}
\usepackage{booktabs}
\usepackage{tikz}
\usepackage{tcolorbox}
\usepackage{multicol}
\usepackage{enumitem}

\tcbuselibrary{skins,breakable}

% ── TODO box ──────────────────────────────────────────────────────────────────
\newtcolorbox{todobox}[1][]{
  colback=yellow!10,
  colframe=orange!80!black,
  fonttitle=\bfseries,
  title={TODO},
  #1
}

% ── Image placeholder box ────────────────────────────────────────────────────
\newtcolorbox{imgplaceholder}[1][]{
  colback=gray!10,
  colframe=gray!60!black,
  fonttitle=\bfseries,
  title={Figure Placeholder},
  #1
}

% ── Code listing style ───────────────────────────────────────────────────────
\definecolor{codebg}{RGB}{248,248,248}
\definecolor{keyword}{RGB}{0,0,180}
\definecolor{string}{RGB}{0,128,0}
\definecolor{comment}{RGB}{128,128,128}
\definecolor{pykey}{RGB}{175,0,219}

\lstdefinelanguage{orca}{
  keywords={model, agent, tool, task, knowledge, workflow, trigger, schema, input, let, edge, true, false, null},
  keywordstyle=\color{keyword}\bfseries,
  commentstyle=\color{comment}\itshape,
  stringstyle=\color{string},
  morecomment=[l]{//},
  morecomment=[l]{\#},
  morestring=[b]",
  morestring=[s]{```}{```},
  sensitive=true,
}

\lstdefinestyle{orcaStyle}{
  language=orca,
  backgroundcolor=\color{codebg},
  basicstyle=\ttfamily\scriptsize,
  breaklines=true,
  frame=single,
  framerule=0.4pt,
  rulecolor=\color{gray!50},
  numbers=left,
  numberstyle=\tiny\color{gray},
  tabsize=2,
  captionpos=b,
  aboveskip=6pt,
  belowskip=4pt,
  xleftmargin=1.5em,
  framexleftmargin=1.5em,
  showspaces=false,
  showstringspaces=false,
}

\lstdefinestyle{pythonStyle}{
  language=Python,
  backgroundcolor=\color{codebg},
  basicstyle=\ttfamily\scriptsize,
  breaklines=true,
  frame=single,
  framerule=0.4pt,
  rulecolor=\color{gray!50},
  numbers=left,
  numberstyle=\tiny\color{gray},
  tabsize=2,
  captionpos=b,
  aboveskip=6pt,
  belowskip=4pt,
  xleftmargin=1.5em,
  framexleftmargin=1.5em,
  keywordstyle=\color{pykey}\bfseries,
  stringstyle=\color{string},
  commentstyle=\color{comment}\itshape,
}

\lstset{style=orcaStyle}

% ── Metadata ─────────────────────────────────────────────────────────────────
\title{Orca: A Declarative Domain-Specific Language for AI Agent Orchestration}

\author{Thakee Nathees}
\email{t.abdulmaj@student.curtin.edu.au}
\affiliation{%
  \institution{Curtin University}
  \department{Department of Computing}
  \city{Perth}
  \country{Australia}
}

\begin{document}

\begin{abstract}
This paper presents \textbf{Orca}, a declarative domain-specific language
(DSL) tailored for the orchestration of autonomous AI agents. Orca is
developed as a robust abstraction over general-purpose host languages such
as Python~\cite{vanrossum2009python} and libraries such as
LangGraph~\cite{langgraph2024}, which form a common programming stack for
LLM-based and agentic systems~\cite{crewai2024,autogen2024}. Orca treats
agent configurations, tool bindings, and execution graphs as first-class
primitives in a declarative notation, enabling compile-time validation of
agent topologies and tool-binding constraints. The Orca compiler implements
a four-stage pipeline (lexer, Pratt parser, semantic analyzer, code
generator) in Go, embeds a self-hosted schema system, and generates
source-mapped Python/LangGraph code. We position Orca against imperative
frameworks and YAML-first configuration approaches, and summarize the
language design and implementation features that support static analysis,
extensibility, and IDE integration. An evaluation across four use cases of
increasing complexity shows that Orca reduces physical source lines of
code by roughly $27\%$ against Python/LangGraph and $44\%$ against
CrewAI, while statically catching four classes of
misconfiguration---reference, type, schema, and workflow
errors---that the imperative baselines surface only at runtime.
\end{abstract}

\maketitle

% ══════════════════════════════════════════════════════════════════════════════
\section{Introduction}
\label{sec:intro}
% ══════════════════════════════════════════════════════════════════════════════

This paper presents Orca, a declarative, domain-specific language tailored
for the orchestration of autonomous AI agents. Orca is a novel programming
language developed as a more robust abstraction to languages such as
Python~\cite{vanrossum2009python} and libraries such as
LangGraph~\cite{langgraph2024}, which form the modern programming stack when
it comes to LLMs and agentic systems~\cite{crewai2024,autogen2024}. It aims
to address the limitations of the current imperative
approach~\cite{docker2025cagent} by introducing an alternative paradigm
where agent configurations, tool bindings and execution graphs are treated
as first-class primitives and represented in a declarative manner. This
shift eliminates a significant amount of boilerplate code necessitated by
the incumbent frameworks and facilitates the lifting of dynamic
misconfiguration errors into the static domain, allowing for compile-time
validation of agent topologies and tool-binding constraints.

Orca adopts an HCL-inspired~\cite{terraform2024} declarative block syntax
where developers simply describe the desired state of the system, and the
Orca compiler translates these declarations into fully functional
Python/LangGraph code, performing static analysis along the way. This paper
makes the following contributions:
\begin{enumerate}[leftmargin=*]
  \item Declarative language design for AI agent orchestration, with a
        uniform block syntax covering models, agents, tools, tasks,
        workflows, and triggers.
  \item Four-stage compiler (lexer, Pratt parser, semantic analyzer, code
        generator) implemented in Go, featuring error-tolerant parsing and a
        unified type system.
  \item Static analysis framework that catches reference errors, type
        mismatches, and schema violations at compile time, with diagnostic
        codes suitable for IDE integration.
  \item Source-mapped code generation that annotates every line of
        generated Python with its originating source location, enabling
        DSL-level debugging.
\end{enumerate}

% ══════════════════════════════════════════════════════════════════════════════
\section{Background}
\label{sec:background}
% ══════════════════════════════════════════════════════════════════════════════

\subsection{The Problem}
\label{sec:background:problem}

\subsubsection{Rise of Agentic Systems}
\label{sec:background:rise}

With the mainstreaming of the LLM (Large Language Model) in
2023~\cite{benaich2023stateofai}, there was a significant shift in the
software systems created~\cite{autogen2024}. There has been a significant
increase in AI-based systems developed, and the need for AI features has
risen tremendously. This has led to rapid developments in the industry,
culminating in more inherently complex systems being built using innovative
LLM technologies. To further accentuate this complexity, there has been a
rising number of inference provisioning methods~\cite{autogen2024}.

When considering the LLM-related technologies that have been developed, it
would be natural to start with AI Agents. These are autonomous software
systems which have the ability to perceive and interact with their
environment~\cite{crewai2024}. Agents can be created by---among other
things---providing an LLM access to a set of ``tools''~\cite{schick2023toolformer}.
These tools may include web search, terminal access, browser automation,
etc.~\cite{schick2023toolformer}. Similarly, there have been other
techniques, such as RAG (Retrieval-Augmented Generation), which attempt to
intelligently manage the LLM context to produce more robust output by
pairing it with ever-evolving knowledge sources such as vector databases
(e.g., Qdrant)~\cite{lewis2020rag}. Further, there have been developments
in techniques such as MCP (Model Context Protocol)~\cite{anthropic2024mcp},
Semantic Routing~\cite{wang2025semanticrouter}, and Agent
Skills~\cite{manias2024skills}. The development of all these technologies has
enabled the creation of many useful AI applications, and they continue to
evolve as time goes on. However, as more such technologies are introduced,
the complexity of the ecosystem continues to trend upwards. This
constitutes the primary gap of the current thesis: a robust abstraction
layer for the orchestration of these technologies.

This need is further amplified by the volume and nature of inference in the
current market. At the time of this writing, LLM inference is provided by
many sources, including OpenAI, Anthropic, Azure, AWS Bedrock, Moonshot AI,
and self-hosted options, just to name a few~\cite{litellm2024}. The array of
provider options has been so wide that there have been many projects to
build singular adapters to multiple providers, such as
LiteLLM~\cite{litellm2024}. This further highlights the need for a robust
AI orchestration abstraction layer.

\subsubsection{AI Agentic Frameworks}
\label{sec:background:frameworks}

Before presenting the case for a further abstraction, the current methods
for the orchestration of agentic systems have to be considered. The key
feature shared among all these technologies is their imperative
nature~\cite{docker2025cagent}.

At the foundational level, many developers initially orchestrate LLMs using
raw Python combined with provider-specific SDKs. This approach requires
engineers to manually wire together API calls, manage execution state,
handle retry logic, and parse unstructured outputs. While this offers
maximum granular control, it rapidly devolves into brittle, unmaintainable
codebase architectures as system complexity scales. Tool binding, context
window management, and memory must be explicitly coded, resulting in vast
amounts of boilerplate simply to establish a basic ReAct (Reasoning and
Acting) loop~\cite{yao2023react}.

To impose structure on this complexity, graph-based frameworks like
LangGraph emerged as an industry standard~\cite{langgraph2024}. LangGraph
treats agent workflows as state machines, explicitly defining nodes
(computational steps) and edges (conditional routing). However, defining
these graphs natively in Python remains highly verbose. The developer is
forced to write extensive imperative code to declare node functions, manage
shared state dictionaries, and compile the graph. Crucially, because the
execution topology is defined at runtime, structural errors---such as
dangling edges, malformed state updates, or incorrect tool schemas---often
fail silently or crash only during active execution.

Parallel to graph-based approaches, multi-agent frameworks such as CrewAI
and Microsoft's AutoGen attempt to simplify orchestration through
role-playing paradigms~\cite{crewai2024,autogen2024}. These libraries allow
developers to define ``personas'' and assign tasks sequentially or
hierarchically. While they mask some of the underlying wiring, they are
still constrained by their imperative host language. Developers must
instantiate classes, pass objects, and handle control flow manually.
Furthermore, these abstractions often obscure the actual execution graph,
making it exceedingly difficult to debug state transitions or trace logic
when an autonomous agent deviates from the expected
path~\cite{crewai2024,autogen2024}.

Ultimately, relying on imperative languages for agent orchestration proves
inefficient and error-prone as these systems evolve. Forcing developers to
manually construct execution graphs and handle low-level wiring obscures the
core intent of the agentic system. The lack of compile-time validation
means that architectural misconfigurations are only caught at runtime,
severely increasing development friction and debugging overhead. This
structural limitation necessitates a shift toward a declarative paradigm,
where the underlying compiler abstracts the ``how'' of execution, allowing
the developer to safely and concisely define the ``what.''

\subsubsection{Domain-Specific Languages}
\label{sec:background:dsl}

To resolve the friction inherent in imperative agent orchestration, the
architectural bottleneck must be addressed at the abstraction layer itself.
This naturally points toward the adoption of a Domain-Specific Language
(DSL) designed explicitly for declarative configuration. A DSL allows
developers to elevate their focus from the mechanics of control
flow---such as manually wiring state dictionaries and handling API retry
logic---to the expression of pure intent. By defining the ``what'' rather
than the ``how,'' a declarative DSL inherently strips away the boilerplate
required to manage execution topologies, shifting the burden of building
the underlying graph to a compiler.

Further, it is interesting to note that this current state of agentic
development closely mirrors the historical evolution of cloud
infrastructure provisioning. In the early days of cloud computing,
engineers relied on imperative scripts to sequentially deploy and configure
resources. These scripts were brittle, difficult to debug, and prone to
state drift, as the developer had to manually account for dependency
resolution and execution order. The paradigm shifted dramatically with the
introduction of Infrastructure as Code (IaC) tools, which allowed
engineers to declare the desired end-state of a system using specialized,
intent-driven languages. The underlying engine then calculated the execution
graph required to achieve that state, fundamentally stabilizing how complex
cloud architectures were deployed.

In the realm of IaC, HashiCorp Configuration Language (HCL), the driving
force behind Terraform, stands out as a highly effective model for this type
of abstraction~\cite{terraform2024}. HCL popularized a block-based syntax
that successfully bridges the gap between human readability and strict
machine validation. It treats infrastructure components (like virtual
machines, subnets, and policies) as first-class primitives. Crucially, it
allows for static validation before execution; resource references, type
constraints, and schema bindings are checked at the compilation stage,
preventing costly dynamic failures during deployment.

Consequently, the ideal solution for the current complexities in AI agent
orchestration could in some way mirror this ``shape.'' The authors hence
concluded that the ecosystem could benefit from an HCL-inspired,
declarative DSL where models, tools, prompts, and execution workflows are
treated as native primitives. Such a language would allow developers to
define agent topologies through concise block structures while a dedicated
compiler handles the translation into the verbose, underlying imperative
frameworks. By enforcing strict semantic analysis---catching type mismatches,
dangling references, and schema violations at compile time---this abstraction
would bring the rigorous, predictable engineering standards of modern
DevOps directly to the orchestration of autonomous AI systems.

\subsection{Similar Works}
\label{sec:background:similar}

\subsubsection{YAML-Based Configurations}
\label{sec:background:yaml}

To bridge the gap between imperative codebase complexity and declarative
intent, several frameworks have introduced YAML-based configuration layers
as a stopgap measure. For instance, modern iterations of multi-agent
frameworks like CrewAI now encourage developers to define agent personas and
task assignments in YAML files, which are parsed at runtime to hydrate
underlying Python objects~\cite{docker2025cagent}. Similarly, open-source
initiatives such as Docker's cagent provide a standalone YAML-driven runtime
for agent execution~\cite{docker2025cagent}, while enterprise solutions
like Microsoft's Semantic Kernel utilize YAML definitions for prompt
templating and I/O schema constraints~\cite{docker2025cagent}.

While the rapid adoption of YAML configurations clearly validates the
industry's demand for declarative orchestration, YAML remains fundamentally
inadequate as a robust engineering abstraction. Because YAML is purely a
data serialization format, it possesses no native semantics for variables,
reference tracking, or type safety. If a YAML-defined agent attempts to
invoke a non-existent tool, or if a variable assignment violates a type
constraint, the YAML parser remains completely agnostic to the error.
Consequently, these structural misconfigurations are entirely invisible
until the application is actively running, leading to costly dynamic
failures. This highlights the critical necessity for a true compiled DSL like
Orca, which replaces runtime parsing with rigorous, compile-time static
analysis.

\subsubsection{LLM Provider Adapters}
\label{sec:background:adapters}

For the sake of a comprehensive review, the role of LLM provider adapters
within the current orchestration stack must also be acknowledged. As the
landscape of inference providers has fragmented across entities like
OpenAI, Anthropic, and various self-hosted open-source models, libraries
such as LiteLLM have emerged to establish a unified
interface~\cite{litellm2024}. These tools allow developers to hot-swap
underlying language models by simply changing a string reference, completely
abstracting the boilerplate required to manage disparate API endpoints,
authentication headers, and network retry logic~\cite{litellm2024}.

The widespread adoption of these adapters demonstrates the undeniable
success and necessity of abstraction layers in reducing LLM integration
friction. However, their scope is strictly constrained to standardizing model
API routing. Provider adapters do not attempt to solve the broader, more
complex challenges of agent orchestration, such as managing execution state,
wiring deterministic tool graphs, or orchestrating multi-agent workflows.
While adapters successfully standardize the isolated ``inference'' node of
a given system, Orca extends this philosophy to abstract and validate the
entire execution topology.

% ══════════════════════════════════════════════════════════════════════════════
\section{Literature Review}
\label{sec:litreview}
% ══════════════════════════════════════════════════════════════════════════════

The field of artificial intelligence is currently undergoing a significant
paradigm shift, transitioning from monolithic large language model (LLM)
interactions to sophisticated, autonomous multi-agent systems (MAS) capable
of executing complex, multi-step tasks~\cite{crewai2024,autogen2024}. This
evolution represents a move beyond passive text generation toward active
participants in problem-solving that perceive their environment, make
decisions, and take actions to achieve specific goals~\cite{crewai2024,autogen2024}.
At the heart of this transition lies the challenge of orchestration---the
management of coordination, communication, and state across diverse agents
and tools~\cite{autogen2024,anthropic2024mcp}. While imperative frameworks
such as LangGraph and multi-agent libraries like CrewAI have provided the
initial foundation for this era, they introduce significant cognitive
overhead, boilerplate code, and reliability risks due to their
runtime-heavy, imperative nature~\cite{docker2025cagent,crewai2024}. This
report evaluates the current state of the agentic stack and establishes the
technical necessity for Orca, a declarative domain-specific language (DSL)
designed to bring the rigorous standards of infrastructure engineering to
the world of AI agents.

\subsection{The Genesis of the Agentic Era and the 2023 LLM Inflection Point}
\label{sec:litreview:genesis}

The mainstreaming of Large Language Models in 2023 served as the primary
catalyst for a fundamental change in software
architecture~\cite{benaich2023stateofai}. According to industry analysis,
the rapid adoption of products like ChatGPT transformed LLMs into one of
the fastest-growing internet products in history, leading more than
two-thirds of organizations to plan increased AI
investments~\cite{benaich2023stateofai}. This surge in interest has
culminated in the development of inherently complex systems that leverage
LLMs not just as information retrieval tools, but as reasoning
engines~\cite{crewai2024}.

As the ecosystem matured throughout 2023 and 2024, the limitations of
standalone models became apparent, particularly regarding factual recall and
real-world grounding~\cite{lewis2020rag}. This led to the rise of
``agentic'' systems---autonomous software entities that utilize LLMs to plan
and execute actions via external tools~\cite{yao2023react,schick2023toolformer}.
The integration of these components has necessitated a robust orchestration
layer capable of managing the ever-increasing complexity of inference
provisioning and model interaction~\cite{autogen2024}.

\begin{table}[t]
  \centering
  \caption{Agentic capabilities and supporting literature.}
  \label{tab:agentic-capabilities}
  \small
  \begin{tabular}{@{}p{0.28\linewidth}p{0.52\linewidth}p{0.14\linewidth}@{}}
    \toprule
    \textbf{Capability} & \textbf{Description} & \textbf{Source} \\
    \midrule
    Reasoning Traces & Internal ``thoughts'' that allow a model to plan before acting & \cite{yao2023react} \\
    Tool Interaction & Ability to invoke APIs, search the web, or execute code & \cite{schick2023toolformer} \\
    Environment Perception & Processing feedback from external actions to update internal state & \cite{crewai2024} \\
    Stateful Persistence & Maintaining conversation and task history over long horizons & \cite{langgraph2024} \\
    \bottomrule
  \end{tabular}
\end{table}

The rise of agentic systems has also been supported by specialized
techniques such as Retrieval-Augmented Generation (RAG), which pairs LLMs
with dynamic knowledge sources like vector databases to reduce hallucinations
and improve contextual accuracy~\cite{lewis2020rag}. However, as the number
of available technologies---including the Model Context Protocol (MCP),
Semantic Routing, and Agent Skills---continues to trend upwards, the primary
gap remains a robust abstraction layer to coordinate these disparate parts
into a stable whole~\cite{anthropic2024mcp,wang2025semanticrouter}.

\subsection{Paradigms of Orchestration: The Limitations of Imperative Frameworks}
\label{sec:litreview:imperative}

The current software stack for LLM applications is heavily reliant on
Python~\cite{vanrossum2009python} and specialized SDKs~\cite{autogen2024}.
Python's dominance in the field stems from its clear syntax and extensive
library ecosystem, which has made it the de-facto language for machine
learning and data science~\cite{vanrossum2009python}. However, when applied
to the orchestration of autonomous agents, the imperative nature of Python
introduces significant challenges~\cite{docker2025cagent,crewai2024}.

\subsubsection{The Burden of Imperative Manual Wiring}
\label{sec:litreview:wiring}

In the initial stages of agent development, engineers often rely on raw
Python combined with provider-specific SDKs to handle API calls, retry
logic, and state management. This approach, while offering granular control,
forces developers to manually handle low-level mechanics that are
peripheral to the core logic of the agent~\cite{docker2025cagent}. The
necessity of explicitly coding tool bindings, context window management, and
memory results in vast amounts of boilerplate code just to establish a basic
Reasoning and Acting (ReAct) loop~\cite{yao2023react}.

Imperative orchestration requires the developer to define exactly how a
system should reach a state, rather than what the state should
be~\cite{docker2025cagent}. This ``fire and forget'' logic is particularly
fragile in the face of transient errors~\cite{docker2025cagent}. If an
external API encounters a brief period of latency or failure, an imperative
script often crashes or enters an inconsistent state, requiring manual
intervention to reconcile processed data against output
tables~\cite{docker2025cagent}.

\subsubsection{Graph-Based Frameworks: LangGraph and Sequential Chains}
\label{sec:litreview:langgraph}

To impose structure on this imperative complexity, graph-based frameworks
like LangGraph~\cite{langgraph2024} have emerged as a standard for building
stateful agents. LangGraph models workflows as state machines, utilizing a
graph-based architecture to manage the intricate relationships between
various nodes (computational steps) and edges (transitions)~\cite{langgraph2024}.

\begin{table}[t]
  \centering
  \caption{LangGraph components and structural constraints.}
  \label{tab:langgraph-components}
  \small
  \begin{tabular}{@{}p{0.18\linewidth}p{0.42\linewidth}p{0.34\linewidth}@{}}
    \toprule
    \textbf{Component} & \textbf{Technical Role} & \textbf{Structural Constraint} \\
    \midrule
    Nodes & Python functions encoding agent logic & Must accept and return state \\
    Edges & Determine the next node based on transitions & Fixed or conditional branching \\
    State & Shared data structure acting as memory & Typically a TypedDict or Pydantic model \\
    Compilation & Validates graph structure at runtime & Checks for orphaned nodes \\
    \bottomrule
  \end{tabular}
\end{table}

Despite its power, LangGraph remains a low-level orchestration
framework~\cite{langgraph2024}. Defining these graphs in native Python is
highly verbose and often obscures the intent of the agentic system.
Crucially, because the execution topology is defined and compiled at
runtime, structural errors such as malformed state updates or incorrect tool
schemas often fail to be detected until active execution, increasing
development friction and debugging overhead~\cite{langgraph2024,docker2025cagent}.

\subsection{Multi-Agent Systems and Role-Based Collaboration}
\label{sec:litreview:mas}

Beyond single-agent orchestration, the industry has embraced multi-agent
systems (MAS) that divide labor among teams of specialized
agents~\cite{crewai2024,autogen2024}. Frameworks like CrewAI~\cite{crewai2024}
and AutoGen~\cite{autogen2024} provide high-level abstractions for this
collaboration.

\subsubsection{AutoGen: Conversational and Adaptive Frameworks}
\label{sec:litreview:autogen}

Microsoft's AutoGen~\cite{autogen2024} introduces a multi-agent conversation
framework where agents interact through automated chats. AutoGen agents are
highly customizable and can operate in various modes, combining LLMs with
tools and human feedback~\cite{autogen2024}. Research indicates that
AutoGen's conversational approach allows agents to exchange information and
opinions to solve complex tasks, often outperforming single-agent
solutions~\cite{autogen2024}. However, the open-ended nature of
conversational agents can make it difficult to trace the actual execution
graph or debug state transitions when an agent deviates from the expected
path~\cite{crewai2024,autogen2024}.

\subsubsection{CrewAI: Hierarchical Orchestration and Roles}
\label{sec:litreview:crewai}

CrewAI~\cite{crewai2024} emphasizes a role-based collaboration paradigm. In
this framework, independent agents are organized into a ``crew,'' where each
member has a specific role (e.g., Researcher, Analyst, Writer) and a
goal~\cite{crewai2024}.

\begin{table}[t]
  \centering
  \caption{CrewAI philosophy and reliability impact.}
  \label{tab:crewai-philosophy}
  \small
  \begin{tabular}{@{}p{0.22\linewidth}p{0.38\linewidth}p{0.34\linewidth}@{}}
    \toprule
    \textbf{Feature} & \textbf{CrewAI Philosophy} & \textbf{Impact on Reliability} \\
    \midrule
    Delegation & Manager agents autonomously assign subtasks & Enables complex, hierarchical work \\
    Processes & Supports sequential and hierarchical execution & Increases predictability \\
    Tools & Extensive library of pre-built tool integrations & Accelerates common development \\
    Memory & Vector database integration for context & Improves long-horizon task coherence \\
    \bottomrule
  \end{tabular}
\end{table}

While CrewAI simplifies orchestration, it remains constrained by the
imperative nature of its host language. Developers must still manually
instantiate classes and handle control flow, and structural
misconfigurations are typically only caught during runtime~\cite{crewai2024}.

\subsection{Technical Necessity for Orca: The Case for Declarative Orchestration}
\label{sec:litreview:necessity}

To resolve the friction inherent in current agent orchestration, the
architectural bottleneck must be addressed at the abstraction layer itself.
The evolution of agentic systems has reached a point where reliance on
imperative scripts is no longer sufficient for production-grade
reliability~\cite{docker2025cagent,manias2024skills}. The following factors
establish the core technical necessity for the Orca language:

\begin{enumerate}[leftmargin=*]
  \item \textbf{Desired-State Management and Reconciliation.} The current
        state of agent development closely mirrors the historical evolution
        of cloud infrastructure~\cite{docker2025cagent}. Early
        cloud engineering relied on imperative scripts that were brittle and
        prone to state drift~\cite{burns2016borg}. The introduction of
        Infrastructure as Code (IaC) tools like Terraform and
        Kubernetes~\cite{burns2016borg} allowed operators to declare a
        desired end-state. Kubernetes solved the infrastructure crisis
        through a ``reconciliation loop''---a continuous control loop that
        monitors actual state, compares it to the declared state, and takes
        corrective action~\cite{burns2016borg,docker2025cagent}. This model is
        essential for AI agents because transient LLM failures or API errors
        are unavoidable; a declarative system like Orca handles retries and
        state recovery as architectural invariants~\cite{docker2025cagent}.

  \item \textbf{Resolving the ``Spec Gap.''} Engineering teams using AI
        agents frequently encounter the ``spec gap''---a drift problem where
        the natural language specification and the actual implementation in
        code begin to diverge~\cite{docker2025cagent}. Because AI generation
        is non-deterministic, vagueness in the specification leads agents to
        fill gaps with assumptions~\cite{docker2025cagent}. By using a
        formal DSL like Orca, the specification is the implementation.
        Declarative outcomes produce more stable agent behavior because they
        describe desired constraints rather than prescribing every step,
        allowing the agent to reason within established boundaries while the
        compiler ensures structural integrity~\cite{docker2025cagent,wang2025semanticrouter}.

  \item \textbf{Compile-Time Static Analysis.} The critical necessity for Orca
        is the ability to perform rigorous compile-time static analysis. In
        imperative frameworks, reference errors, type mismatches, and schema
        violations are invisible until the application crashes during
        execution~\cite{docker2025cagent,terraform2024}. Orca's compiler
        facilitates the lifting of these dynamic misconfiguration errors into
        the static domain, catching:
        \begin{itemize}[leftmargin=*]
          \item \textbf{Reference Errors:} Ensuring all agents, tools, and
                workflows referenced in a manifest exist~\cite{terraform2024}.
          \item \textbf{Type Mismatches:} Validating that I/O schemas between
                agents and tools are compatible before generation~\cite{schick2023toolformer}.
          \item \textbf{Schema Violations:} Checking that block configurations
                adhere to the expected structure~\cite{terraform2024}.
        \end{itemize}
        This shift toward compile-time validation brings the predictable
        engineering standards of DevOps~\cite{terraform2024} directly to the
        orchestration of autonomous AI systems.
\end{enumerate}

\subsection{Comparative Review of Stopgap Measures: YAML and Adapters}
\label{sec:litreview:stopgap}

The industry's demand for declarative orchestration has led to several
stopgap measures, most notably YAML-based configuration layers and LLM
provider adapters~\cite{litellm2024,docker2025cagent}.

\subsubsection{YAML-Based Configurations: Data without Semantics}
\label{sec:litreview:yaml-semantics}

Frameworks like Docker's cagent~\cite{docker2025cagent} and Semantic Kernel
utilize YAML to decouple agent personas and task assignments from
underlying code. This configuration-first philosophy is designed for
portability and execution speed~\cite{docker2025cagent}. However, YAML is a
data serialization format, not a programming language~\cite{terraform2024}. It
possesses no native semantics for variables or reference tracking. A
YAML-defined agent misconfiguration remains invisible until runtime, leading
to costly dynamic failures~\cite{docker2025cagent}. Orca replaces runtime
parsing with a dedicated compiler that understands the semantic relationships
between agents and their environment.

\subsubsection{LLM Provider Adapters: Standardizing the Inference Node}
\label{sec:litreview:litellm}

The role of LLM provider adapters, such as LiteLLM~\cite{litellm2024},
demonstrates the success of abstraction layers. LiteLLM provides a unified
interface to call 100+ LLMs, abstracting authentication, API endpoints, and
retry logic~\cite{litellm2024}. While adapters successfully standardize the
``inference'' node, they do not solve the broader challenges of agent
coordination or execution state~\cite{litellm2024}. Orca extends this
standardization to the entire execution topology.

\subsection{Standardized Protocols and Reusable Agentic Primitives}
\label{sec:litreview:protocols}

The development of Orca occurs alongside the broader movement toward
standardization, particularly the Model Context Protocol
(MCP)~\cite{anthropic2024mcp} and ``Agent Skills''~\cite{crewai2024}.

\subsubsection{Model Context Protocol (MCP): The Universal Interface}
\label{sec:litreview:mcp}

Introduced by Anthropic in November 2024, the Model Context Protocol (MCP)
is an open standard designed to bridge the gap between AI assistants and
the world of data and tools~\cite{anthropic2024mcp}. MCP provides a universal
interface for reading files and executing functions, acting as a ``USB-C for
AI''~\cite{anthropic2024mcp}. For an orchestration language like Orca, MCP
provides a standardized ``tool binding'' primitive, allowing the compiler
to validate I/O schemas between models and tools at compile
time~\cite{anthropic2024mcp,docker2025cagent}.

\subsubsection{Agentic Skills: Modular Procedural Capabilities}
\label{sec:litreview:skills}

``Agentic Skills'' are reusable, callable modules that package procedural
knowledge with explicit applicability conditions and execution
policies~\cite{manias2024skills}. Skills differ from tools in that they
package multi-step sequences and explicit termination
criteria~\cite{manias2024skills}. Orca treats these skills as first-class
primitives, allowing developers to define a uniform block syntax for agents
that includes specific expertise and activation triggers~\cite{manias2024skills}.

\subsection{Synthesis and Conclusion}
\label{sec:litreview:synthesis}

The transition toward autonomous multi-agent systems represents the next
frontier of AI application development, but current imperative orchestration
frameworks create a ``fragility gap''~\cite{docker2025cagent,manias2024skills}.
Orca addresses these structural limitations by introducing a declarative
paradigm where agent configurations and workflows are first-class
primitives. By adopting an HCL-inspired syntax~\cite{terraform2024} and a
compiler with rigorous static analysis, Orca provides operational stability,
structural integrity, and developer productivity~\cite{terraform2024,docker2025cagent}.
Orca brings the engineering standards of modern DevOps to agentic systems,
ensuring that the next generation of AI applications is as stable and
maintainable as the infrastructure they run on~\cite{burns2016borg,terraform2024}.

% ══════════════════════════════════════════════════════════════════════════════
\section{Design}
\label{sec:design}
% ══════════════════════════════════════════════════════════════════════════════

Orca is a declarative domain-specific language designed to abstract the
structural complexity of AI agent orchestration. Rather than requiring
developers to manually write procedural execution loops, state management,
and framework-specific wiring, Orca allows developers to define the desired
topology of an agentic system. The compiler translates these declarations
into fully functional, imperative framework code (currently targeting
LangGraph).

\subsection{Design Principles}
\label{sec:design:principles}

The architecture of Orca is guided by four core principles designed to
address the inherent limitations of imperative orchestration:
\begin{enumerate}[leftmargin=*]
  \item \textbf{Declarative Intent:} The language focuses exclusively on the
        ``what'' rather than the ``how.'' By eliminating general-purpose
        control flow loops (\texttt{while}, \texttt{for}) and conditional
        branching outside of explicit graph definitions, developers are
        constrained to describing state and relationships, producing highly
        predictable architectures.
  \item \textbf{Static Safety:} The reliance on late-failing,
        runtime-heavy dictionaries in Python is replaced with compile-time
        validation. Reference errors, missing required fields, and type
        mismatches are caught before execution.
  \item \textbf{Framework Independence:} While modern orchestration is
        tightly coupled to specific SDKs (e.g., LangGraph, AutoGen), Orca
        separates the system definition from the execution environment. The DSL
        serves as a universal frontend, enabling future compiler backends to
        target different execution frameworks without altering the source
        configuration.
  \item \textbf{Debuggability:} To ensure the abstraction does not obscure
        underlying execution, Orca is designed with source-mapped code
        generation, allowing runtime behaviors to be traced directly back to
        their originating \texttt{.oc} declarations.
\end{enumerate}

\subsection{Syntax and Program Structure}
\label{sec:design:syntax}

Orca adopts an HCL-inspired block syntax. An Orca program (an \texttt{.oc}
file) consists of a flat sequence of named blocks. There are no imports, no
top-level expressions, and no standalone function definitions. Every
component of the system is encapsulated within a block consisting of a
keyword, a unique identifier, and a brace-delimited body of assignments.

\begin{lstlisting}[caption={A standard Orca block definition.},label={lst:orca-block}]
// A standard Orca block definition
model gpt4 {
  provider    = "openai"
  model_name  = "gpt-4o"
  temperature = 0.2
}
\end{lstlisting}

The language supports standard literals (strings, integers, floats,
booleans, and null) as well as parameterized collections (lists and maps).
A notable syntactic feature is the implementation of multi-line raw strings
(delimited by triple backticks), which feature automatic indentation
normalization. The compiler utilizes the closing delimiter's column position
as the baseline, stripping leading whitespace to prevent indentation from
leaking into system prompts or inline Python code.

\subsection{First-Class Orchestration Primitives}
\label{sec:design:primitives}

Orca elevates the conceptual components of multi-agent systems into native
language primitives. There are eight implemented block types, with four
serving as the primary orchestration mechanisms:
\begin{itemize}[leftmargin=*]
  \item \textbf{Models and Agents:} The \texttt{model} block configures LLM
        provider settings (e.g., OpenAI, Anthropic). The \texttt{agent}
        block defines an autonomous entity, binding a specific model to a
        system prompt (persona) and a list of callable tools.
  \item \textbf{Tools:} The \texttt{tool} block defines external
        integrations. Orca supports both dotted Python import paths for
        existing libraries and inline Python function definitions (using raw
        string language tags, e.g., \texttt{py}), which the compiler
        automatically extracts and injects into the generated environment.
  \item \textbf{Workflows:} Execution graphs are defined using the
        \texttt{workflow} block. Rather than utilizing assignments, the body
        of a workflow consists of bare expressions utilizing the arrow
        operator (\texttt{->}). This provides a native, highly readable
        syntax for defining execution topologies:
\end{itemize}

\begin{lstlisting}[caption={Workflow topology using the arrow operator.},label={lst:orca-workflow}]
workflow content_pipeline {
  researcher -> writer -> reviewer
}
\end{lstlisting}

\noindent This single line of Orca naturally transpiles to the complex
node-and-edge state machine definitions required by frameworks like
LangGraph.

\subsection{The Structural Type System and Schemas}
\label{sec:design:types}

To provide strict static analysis, Orca utilizes a structural type system
built around four type kinds: \texttt{BlockRef} (named primitive types and
block references), parameterized \texttt{List}, parameterized \texttt{Map},
and \texttt{Union} types (represented by the pipe operator \texttt{|}).

The contract between the developer and the compiler is governed by block
schemas. Every block type in Orca possesses a rigid schema defining its
valid fields, required status, and expected types. For example, the
\texttt{agent} schema strictly requires a model reference and a persona
string, while \texttt{tools} remains an optional list.

A critical architectural decision in Orca's design is a self-bootstrapped
type system. The built-in schemas defining the core language constructs
(such as \texttt{model} and \texttt{agent}) are not hardcoded into the Go
compiler's source; rather, they are defined in native Orca syntax within an
embedded \texttt{builtins.oc} file. This self-hosting approach unifies
primitives and user-defined structures under the \texttt{BlockRef} type,
resulting in a highly extensible type checker. Developers can define their
own custom data structures using the \texttt{schema} block, which the
compiler validates seamlessly alongside standard language primitives.

\subsection{Extensibility and Developer Experience}
\label{sec:design:dx}

Orca incorporates several features to ensure the language is suitable for
complex, production-grade engineering:
\begin{itemize}[leftmargin=*]
  \item \textbf{Variables and Indirection:} The \texttt{let} block allows
        developers to define named constants and environment configurations.
        The compiler performs static constant folding, traversing \texttt{let}
        indirections at compile time to resolve underlying values (e.g.,
        dynamically resolving LLM provider strings for code generation).
  \item \textbf{Runtime Inputs:} The \texttt{input} block declares typed
        parameters required at runtime, preventing hardcoded secrets and
        facilitating environment-specific deployments.
  \item \textbf{Annotations:} Blocks and fields can be decorated with
        metadata. The \texttt{@desc} annotation attaches documentation for IDE
        hover integration, while the \texttt{@suppress} directive allows
        developers to selectively bypass specific compiler diagnostics (e.g.,
        \texttt{@suppress("unknown-field")}), providing necessary flexibility
        during the active development of new system features.
\end{itemize}

% ══════════════════════════════════════════════════════════════════════════════
\section{Implementation}
\label{sec:implementation}
% ══════════════════════════════════════════════════════════════════════════════

This section details the practical realization of the Orca language and its
underlying compiler. The implementation is broken down into the core
pipeline required to transpile Orca to Python/LangGraph, followed by the
advanced architectural enhancements that elevate the language into a
production-ready developer tool.

\subsection{Basic Implementation}
\label{sec:implementation:basic}

The basic implementation of Orca encompasses the four-stage compiler
pipeline necessary to convert a declarative \texttt{.oc} source file into an
imperative, executable LangGraph state machine. The compiler is entirely
written in Go and operates via a linear, strictly typed pipeline:
\texttt{Lexer} $\rightarrow$ \texttt{Parser} $\rightarrow$ \texttt{Semantic
Analyzer} $\rightarrow$ \texttt{Code Generator}.

\paragraph{Features and operation.} The system operates by passing source
code through distinct transformation phases without relying on a separate
Intermediate Representation (IR).
\begin{enumerate}[leftmargin=*]
  \item \textbf{Lexical Analysis:} The scanner processes the raw source text
        into a stream of typed tokens, mapping precise source coordinates
        (line, column, end-line, end-column) to enable strict downstream
        diagnostic reporting.
  \item \textbf{Parsing:} The parser constructs the Abstract Syntax Tree
        (AST). It utilizes a hybrid approach: recursive descent for
        top-level block structures and a Pratt parser to handle expression
        precedence (such as the \texttt{->} workflow edge operator).
  \item \textbf{Semantic Analysis:} This phase validates the AST against
        Orca's structural type system. It operates in three passes: symbol
        table construction, schema registration, and per-block validation
        (checking for type mismatches, missing fields, and undefined
        references).
  \item \textbf{Code Generation:} The final validated AST is traversed by
        the LangGraph backend to emit Python code. Blocks are translated into
        LangGraph constructs, with provider strings mapping to LangChain
        imports and workflow edges mapping to StateGraph edge definitions.
\end{enumerate}

\paragraph{Data structures.} The implementation relies heavily on a unified
interface-driven AST. The central data structure is the \texttt{BlockBody},
which encapsulates both top-level statements (\texttt{model}, \texttt{agent})
and inline expressions (\texttt{schema \{ \ldots\ \}}). By sharing this
structure, the semantic analyzer and code generator can traverse and validate
both top-level and inline blocks through the same code paths. Additionally,
the \texttt{SymbolTable} data structure maps identifiers to their resolved
\texttt{BlockRef} types, enabling robust cross-referencing.

\paragraph{Implementation challenges.} The most significant challenge during
the basic implementation was the development of the error-tolerant parser.
Standard parsers halt on syntax errors, but Orca is designed to support
real-time IDE integration. Implementing recovery
strategies---specifically the \texttt{syncToBlockEnd} and
\texttt{syncToNextAssignment} algorithms---was complex. These functions
allow the parser to skip malformed tokens and generate a partial AST,
ensuring the semantic analyzer can still provide diagnostics for the rest of
the file without entering an infinite loop. Additionally, allowing block
keywords (like \texttt{model} or \texttt{tool}) to also act as valid
assignment identifiers required careful lookahead logic in the parser to
avoid syntax ambiguity.

\begin{imgplaceholder}[title={Screenshot: successful \texttt{orca build}}]
Terminal output showing the \texttt{orca build} command successfully
generating the \texttt{build/} directory with source-mapped Python output.
\end{imgplaceholder}

\subsection{Self-Bootstrapped Schema System}
\label{sec:implementation:bootstrap}

A characteristic design decision of the Orca compiler is that the schemas
describing built-in block types---\texttt{model}, \texttt{agent},
\texttt{tool}, \texttt{workflow}, \texttt{input}, and the primitive types
\texttt{str}, \texttt{int}, \texttt{float}, \texttt{bool}, \texttt{list},
\texttt{map}, \texttt{any}, \texttt{null}---are not hardcoded as Go
structures inside the compiler. Instead, they are written in ordinary Orca
syntax in a file named \texttt{builtins.oc} that is embedded into the
compiler binary with Go's \texttt{//go:embed} directive and parsed by the
same lexer, parser, and analyzer that process user code. Listing~\ref{lst:builtins}
reproduces a fragment of this file: the \texttt{model} schema declares its
fields and their types using exactly the notation available to end users.

\begin{lstlisting}[caption={Excerpt from the embedded \texttt{builtins.oc}: the \texttt{model} schema is declared in the same language the compiler compiles.},label={lst:builtins}]
schema str   {}
schema int   {}
schema float {}
schema bool  {}

schema model {
  provider    = str
  model_name  = str | model
  api_key     = str | null
  base_url    = str | null
  temperature = float | null
}

schema agent {
  model         = str | model
  persona       = str
  tools         = list[tool] | null
  output_schema = schema | null
}
\end{lstlisting}

Three properties follow from this self-hosted arrangement. First, primitive
types and user-defined schemas are represented uniformly as instances of the
same \texttt{BlockRef} type kind in the analyzer's symbol table; there is no
two-level type system that distinguishes ``built-in'' from ``user'' types.
Second, a user's \texttt{schema Ticket \{ \ldots\ \}} block participates in
the same type-checking pipeline as \texttt{str} or \texttt{agent}, which
makes custom structured I/O first-class rather than bolted on. Third,
evolving the language---adding an optional field to \texttt{agent} or
introducing a new block kind---is an edit to a single \texttt{.oc} file
rather than to analyzer Go code, which keeps the type system discoverable
from within the language it describes. The \texttt{@suppress("duplicate-block")}
annotations visible in Listing~\ref{lst:builtins} illustrate a secondary
benefit: the same diagnostic-suppression mechanism exposed to users is what
the language uses to bootstrap itself.

\subsection{The Pratt Parser and the Arrow Operator}
\label{sec:implementation:pratt}

Orca's expression grammar is parsed with a Pratt (top-down operator
precedence) parser~\cite{kladov2020pratt}, embedded inside a recursive
descent driver that handles block-level structure. Each token carries a
binding power, declared in a compact table with entries for
\texttt{PrecArrow} (\texttt{->}), \texttt{PrecPipe} (\texttt{|}),
\texttt{PrecSum} (\texttt{+ -}), \texttt{PrecProduct} (\texttt{* /}), and
\texttt{PrecAccess} (\texttt{.}, \texttt{[}, \texttt{(}). The parser
consumes a primary expression on the left and then repeatedly folds in
higher-precedence binary operators, producing a left-associative expression
tree without mutual-recursion between per-operator parse functions.

This design earned its place because of three Orca-specific requirements.
\emph{First}, the workflow edge operator \texttt{->} had to be a real
infix expression rather than a specialized sub-grammar, so that
\texttt{researcher -> writer -> reviewer} would parse as a binary
expression tree and later lower to a chain of LangGraph \texttt{add\_edge}
calls in the code generator. \emph{Second}, the union type operator
\texttt{|} (used pervasively in \texttt{builtins.oc}---for example in
\texttt{temperature = float | null}) had to share the same expression
machinery as \texttt{->}, so type expressions and value expressions could
reuse a single parser rather than diverging into two grammars.
\emph{Third}, member access (\texttt{builtins.web\_search}) and
subscripting (\texttt{list[tool]}) needed to bind more tightly than any
binary operator; making them left-denotation entries in the Pratt table
gave them the correct precedence without special casing.

The same framework supports \emph{error-tolerant parsing}. The parser
records a \texttt{Diagnostic} at an offending token and then invokes one of
two recovery routines---\texttt{syncToBlockEnd}, which advances to the next
matching right brace, and \texttt{syncToNextAssignment}, which advances to
the next identifier-equals pattern---so a single malformed field does not
discard the rest of the file. The analyzer can therefore still validate the
well-formed portion, and the Language Server Protocol integration described
in Section~\ref{sec:implementation:enhancements} can surface real-time
diagnostics after every keystroke rather than only on complete documents.

\subsection{Enhancements}
\label{sec:implementation:enhancements}

Beyond the mandatory transpilation pipeline, several advanced features were
implemented to provide a superior developer experience and stabilize the
architecture.
\begin{itemize}[leftmargin=*]
  \item \textbf{Embedded Zero-Dependency Runtime:} Rather than requiring
        developers to install an external Python package via pip, the
        runtime support library (\texttt{orca.py}) is embedded directly into
        the Go compiler binary using the \texttt{//go:embed} directive.
        During code generation, this file is written to the output directory,
        ensuring the generated code is completely self-contained.
  \item \textbf{Compile-Time Constant Folding:} The analyzer includes a
        constant folder that evaluates expressions at compile time. This is
        critical for provider resolution; if a \texttt{model} block
        dynamically references a provider string through a \texttt{let}
        variable, the constant folder traces the indirection so the code
        generator knows exactly which LangChain library to import.
  \item \textbf{Language Server Protocol (LSP) Integration:} The compiler's
        partial-AST tolerance enables a fully functional LSP server built on
        glsp. The server re-parses the document on every keystroke,
        providing real-time diagnostics, hover documentation (reading from
        \texttt{@desc} annotations), and cross-file go-to-definition
        resolution.
\end{itemize}

\begin{imgplaceholder}[title={Screenshot: Orca in VS Code}]
VS Code editor demonstrating an Orca file with hover documentation and
real-time semantic error highlighting provided by the LSP.
\end{imgplaceholder}

\subsection{Orca Studio: Visual Editing on the Canvas}
\label{sec:implementation:studio}

Textual source is the authoritative representation of an Orca program, but
for certain workflows---particularly early design of multi-agent topologies
or explaining a system to non-engineer stakeholders---a visual surface is
more appropriate. \emph{Orca Studio} is a browser-based companion to the
command-line compiler that presents the same \texttt{.oc} source as an
interactive canvas of nodes and edges. It is implemented as a Next.js
application built on the React Flow library~\cite{reactflow}, with the
canvas, inspector, and code-editor panels driven by a shared client-side
store. Models, agents, tools, and workflow steps appear as typed nodes
whose handle colours mirror the structural type they expose; workflow
arrows become React Flow edges whose endpoints are validated against the
same schema as the textual \texttt{->} operator.

Studio is designed around a bidirectional round-trip against the textual
source. When the user drags a node, edits a field in the inspector, or
connects two handles, the change is serialized back into canonical
\texttt{.oc} syntax and reflected in an embedded code editor; conversely,
pasting new \texttt{.oc} source into the editor re-populates the canvas.
The user interface therefore operates as a view over the DSL rather than a
parallel language, and users can move between representations without
losing fidelity. Because the serialization boundary is a syntactic surface
rather than a runtime structure, Studio inherits the static safety of the
compiler: any edit that would produce a reference error, schema violation,
or type mismatch surfaces as a diagnostic in the inspector before the user
invokes any backend.

\begin{imgplaceholder}[title={Screenshot: Orca Studio canvas}]
The Studio canvas displaying a two-agent workflow. Models, agents, and
tools appear as typed nodes; the workflow \texttt{->} operator is rendered
as a directed edge between agent handles.
\end{imgplaceholder}

\begin{imgplaceholder}[title={Screenshot: Studio inspector and code round-trip}]
The inspector panel beside the embedded code editor, showing a field edit
reflected both in the canvas node and in the serialized \texttt{.oc}
source on the right.
\end{imgplaceholder}

\subsection{Commenting Your Code}
\label{sec:implementation:comments}

The source code of the Orca compiler is comprehensively documented in
accordance with standard project guidelines. Because the implementation
language is Go, standard godoc formatting conventions are utilized. Every
package (e.g., \texttt{lexer}, \texttt{parser}, \texttt{analyzer},
\texttt{codegen}) contains package-level documentation explaining its role
in the compilation pipeline. Furthermore, all exported structs, interfaces
(such as the \texttt{Node} interface), and public methods include block
comments detailing their behavior, parameters, and return states. This
documentation is automatically extracted into a searchable HTML format
using Go's native documentation tooling, ensuring the architecture is easily
navigable for inspection and future maintenance.

\subsection{Coding Conventions}
\label{sec:implementation:conventions}

To ensure high code legibility, fewer bugs, and overall structural quality,
the Orca compiler strictly adheres to standard Go idiomatic conventions.
\begin{itemize}[leftmargin=*]
  \item All source code is formatted automatically using the \texttt{gofmt}
        tool prior to compilation, ensuring uniform indentation and spacing.
  \item The codebase enforces strict package isolation. There are no cyclic
        dependencies between the four pipeline stages; data flows strictly
        sequentially from lexing to code generation.
  \item Interfaces are used to cleanly separate statement nodes from
        expression nodes in the AST, leveraging Go's type system to prevent
        illegal AST constructions at the implementation level.
  \item Error handling follows Go conventions, where errors are returned
        explicitly alongside data, rather than relying on exceptions. This
        makes the error-recovery mechanisms in the parser highly explicit and
        traceable.
\end{itemize}

% ══════════════════════════════════════════════════════════════════════════════
\section{Evaluation}
\label{sec:evaluation}
% ══════════════════════════════════════════════════════════════════════════════

The central empirical claims of Orca are that the declarative surface
reduces the authoring effort required to express an agentic system and
that its compiler surfaces misconfigurations earlier than comparable
imperative or data-only approaches. This section evaluates both claims
using four use cases of increasing structural complexity, compared across
four stacks, and an analysis of the compile-time error classes that Orca
catches.

\subsection{Methodology}
\label{sec:evaluation:methodology}

\paragraph{Research question.} \emph{Does the declarative DSL reduce
authoring effort and surface misconfigurations earlier than imperative
baselines?}

\paragraph{Baselines.} Orca is compared against three representative stacks
drawn from the orchestration landscape surveyed in
Section~\ref{sec:litreview}: Python with
LangGraph~\cite{langgraph2024}, which is the imperative target that Orca
itself currently compiles to; Python with CrewAI~\cite{crewai2024}, a
role-based multi-agent framework; and Docker
cagent~\cite{docker2025cagent}, a YAML-driven declarative runtime.
Together these span the three design points that Orca positions itself
against: low-level graph orchestration, high-level role-based
orchestration, and format-level declarative configuration.

\paragraph{Metrics.} Three complementary metrics are reported:
\begin{itemize}[leftmargin=*]
  \item \textbf{Physical source lines of code} (SLOC), counted after
        stripping blank lines and pure-comment lines, using a single awk
        filter applied uniformly to all four variants of every use case so
        that comment density and blank-line style cannot distort the
        comparison.
  \item \textbf{Boilerplate share}, defined as the subset of SLOC whose
        purpose is framework wiring rather than domain description:
        imports, state-dictionary declarations, graph builder calls, and
        scheduler setup. This captures the portion of the program that a
        declarative compiler should be able to generate.
  \item \textbf{Error-class coverage}, a qualitative measure of which
        categories of misconfiguration each stack rejects at author time
        versus first execution.
\end{itemize}

\paragraph{Artefacts.} Each use case is implemented in all four stacks.
The source files live under \texttt{paper/dsl/examples/} and are the
direct inputs to the SLOC measurement, so the numbers below are
reproducible by re-running the line-counting script on the same tree.
Implementations are deliberately written to be idiomatic rather than
minimal; shortcuts that are not representative of production code
(collapsing multi-line definitions, omitting required imports) are
avoided.

\paragraph{Limitations taken up front.} SLOC is a proxy for authoring
effort, not a direct measure of it. It ignores cognitive load, the cost
of discovering framework conventions, and the long-term compounding
effect of refactoring imperative wiring. It also does not capture
readability, which the side-by-side examples in the next subsection are
intended to make qualitatively visible. Section~\ref{sec:evaluation:validity}
returns to these threats.

\subsection{Use Cases}
\label{sec:evaluation:usecases}

Four use cases of increasing structural complexity are evaluated. Each
exercises a progressively larger slice of the language: single agent,
sequential multi-agent, tool use with a scheduled trigger, and finally a
multi-agent pipeline carrying a user-defined typed payload.

\begin{description}[leftmargin=*,style=nextline]
  \item[UC1 Single-agent assistant.] A single model bound to a single
        agent with a short persona. Establishes the lower bound for every
        stack---if Orca does not reduce SLOC here, no more complex case
        will reveal a genuine saving.
  \item[UC2 Two-agent sequential pipeline.] A \emph{researcher} agent
        equipped with a web-search tool feeding a \emph{writer} agent.
        This is the workflow shown on the project landing page and is the
        canonical scenario in which Orca's \texttt{->} operator is
        expected to replace explicit LangGraph edge construction.
  \item[UC3 Scheduled tool-using agent.] A single agent invoking a custom
        Python tool, wired to a cron schedule. This case exercises the
        \texttt{tool} block alongside Orca's first-class \texttt{cron}
        trigger, which neither LangGraph nor CrewAI provide natively; YAML
        cagent likewise requires an external scheduler.
  \item[UC4 Typed multi-agent pipeline.] Three agents---triager,
        responder, quality-assurance---pipelined through a
        \texttt{support\_pipeline} workflow and exchanging a user-defined
        \texttt{Resolution} record. Exercises user-defined \texttt{schema}
        blocks, \texttt{output\_schema} binding on an agent, and the
        typed \texttt{input} mechanism.
\end{description}

For brevity, Figure~\ref{fig:uc2-sidebyside} reproduces only the Orca and
LangGraph implementations of UC2 side by side; the remaining twelve
variants are included in the replication package.

\begin{figure*}[t]
\begin{minipage}[t]{0.46\linewidth}
\begin{lstlisting}[caption={UC2 in Orca (16 SLOC).},label={lst:uc2-orca}]
model claude {
  provider   = "anthropic"
  model_name = "claude-opus-4.6"
}

agent researcher {
  model   = claude
  tools   = [builtins.web_search]
  persona = "You research tech trends."
}

agent writer {
  model   = claude
  persona = "You write concise reports."
}

workflow research_and_write {
  researcher -> writer
}
\end{lstlisting}
\end{minipage}
\hfill
\begin{minipage}[t]{0.50\linewidth}
\begin{lstlisting}[style=pythonStyle,caption={UC2 in Python + LangGraph (22 SLOC).},label={lst:uc2-langgraph}]
from langchain_anthropic import ChatAnthropic
from langchain_community.tools import TavilySearchResults
from langchain_core.messages import SystemMessage
from langgraph.graph import StateGraph, MessagesState

claude = ChatAnthropic(model="claude-opus-4.6")
search = TavilySearchResults()
claude_with_tools = claude.bind_tools([search])

def researcher(state: MessagesState):
    sys = "You research tech trends."
    msgs = [SystemMessage(content=sys)] + state["messages"]
    return {"messages": [claude_with_tools.invoke(msgs)]}

def writer(state: MessagesState):
    sys = "You write concise reports."
    msgs = [SystemMessage(content=sys)] + state["messages"]
    return {"messages": [claude.invoke(msgs)]}

graph = StateGraph(MessagesState)
graph.add_node("researcher", researcher)
graph.add_node("writer", writer)
graph.add_edge("__start__", "researcher")
graph.add_edge("researcher", "writer")
graph.add_edge("writer", "__end__")
app = graph.compile()
\end{lstlisting}
\end{minipage}
\caption{Side-by-side comparison for UC2. The Orca version on the left
expresses the same wiring as the Python version on the right while
eliminating all state, import, and graph-construction boilerplate.}
\label{fig:uc2-sidebyside}
\end{figure*}

\subsection{Lines-of-Code Results}
\label{sec:evaluation:loc}

Table~\ref{tab:loc-results} reports the SLOC measured for each use case
across all four stacks. Numbers are obtained with a single awk filter
(\texttt{NF \&\& \$0 !\textasciitilde\ /\textasciicircum [[:space:]]*(\# \textbar\textbar\ //)/})
applied to the example sources.

\begin{table}[t]
  \centering
  \caption{Physical SLOC per use case and stack. Lower is better.
  Percentages in the final row are the mean reduction achieved by Orca
  relative to the baseline in each column.}
  \label{tab:loc-results}
  \small
  \begin{tabular}{@{}lrrrr@{}}
    \toprule
    \textbf{Use case} & \textbf{Orca} & \textbf{LangGraph} & \textbf{CrewAI} & \textbf{cagent} \\
    \midrule
    UC1 Assistant              & 8  & 13  & 19  & 9  \\
    UC2 Research/Writer        & 16 & 22  & 34  & 21 \\
    UC3 Scheduled tool         & 20 & 27  & 33  & 19 \\
    UC4 Typed multi-agent      & 35 & 47  & 55  & 28 \\
    \midrule
    \textbf{Total}             & \textbf{79} & \textbf{109} & \textbf{141} & \textbf{77} \\
    \textit{Mean reduction vs.\ Orca} & --- & $27.5\%$ & $44.0\%$ & $-2.6\%$ \\
    \bottomrule
  \end{tabular}
\end{table}

Three observations follow from the data. First, Orca produces the smallest
or second-smallest program in every use case, and the gap widens as the
structural complexity grows: UC4 shows a $25.5\%$ reduction against
LangGraph and $36.4\%$ against CrewAI, compared to $38.5\%$/$57.9\%$ in
UC1. The super-linear behaviour reflects that imperative stacks pay a
fixed per-node wiring cost (state dictionaries, add-node calls,
add-edge calls) that Orca collapses into a single \texttt{->} expression
regardless of workflow length.

Second, the comparison against Docker cagent YAML deserves separate
treatment. YAML is within a few lines of Orca for every use case, and is
slightly shorter in aggregate ($77$ vs.\ $79$ SLOC), because both stacks
share the declarative philosophy that eliminates imports and graph-building
calls. The two-line deficit is more than explained by Orca's inline
\texttt{let} and \texttt{@desc} conveniences. What the LOC metric does not
show is that the YAML version of UC4 expresses the \texttt{Resolution}
schema as an inline JSON-schema fragment with no static validation, and
that all four YAML files accept an undefined \texttt{tool} reference or a
misspelled field name without complaint until the runtime loads them. The
next subsection quantifies this difference.

Third, CrewAI's overhead is structural rather than incidental: the
role-based abstraction adds a \texttt{Task} object per agent and requires
an explicit \texttt{context} parameter to wire sequential data flow.
This explains the consistent $40$--$60\%$ SLOC premium over Orca across
use cases.

\subsection{Compile-Error Taxonomy}
\label{sec:evaluation:taxonomy}

The second half of the evaluation concerns \emph{when} errors are caught.
Orca's analyzer emits structured \texttt{Diagnostic} values with a short
machine-readable code, a severity, and a source position (the full
catalogue is given in Appendix~\ref{sec:appendix:diagnostics}). The
diagnostics partition naturally into four buckets; this subsection gives a
minimal broken \texttt{.oc} example from each bucket together with the
Orca diagnostic and the observable failure mode in the imperative
equivalents.

\paragraph{Bucket 1: reference errors.} A workflow or field names an
identifier that is not bound in the symbol table.

\begin{lstlisting}[caption={Reference error: \texttt{writer} is never declared.},label={lst:err-ref}]
agent researcher { model = gpt4 persona = "..." }
workflow pipeline {
  researcher -> writer
}
\end{lstlisting}

\noindent Orca reports \texttt{[undefined-ref] unknown identifier 'writer'}
at parse time. The equivalent LangGraph program
(\texttt{graph.add\_edge("researcher", "writer")}) compiles cleanly and
only fails on the first \texttt{app.invoke}, surfacing as a KeyError
during message routing.

\paragraph{Bucket 2: type mismatches.} A field is assigned a value whose
type is incompatible with the schema entry for that field.

\begin{lstlisting}[caption={Type mismatch: \texttt{temperature} requires \texttt{float}.},label={lst:err-type}]
model gpt4 {
  provider    = "openai"
  temperature = "hot"
}
\end{lstlisting}

\noindent Orca reports \texttt{[type-mismatch] expected float|null, got
str}. The LangGraph equivalent (\texttt{ChatOpenAI(temperature="hot")})
raises a Pydantic \texttt{ValidationError} only when the model client is
instantiated at runtime, inside framework code.

\paragraph{Bucket 3: schema violations.} A required field is absent, an
unknown field is present, or a block is duplicated.

\begin{lstlisting}[caption={Schema violation: required field \texttt{persona} missing.},label={lst:err-schema}]
agent assistant {
  model = gpt4
}
\end{lstlisting}

\noindent Orca reports \texttt{[missing-field] agent requires field
'persona'}. LangGraph's equivalent omission is silent---the node function
is defined without a system prompt---and surfaces only as empty or
malformed LLM output, which is difficult to attribute to a
misconfiguration. cagent YAML accepts an unknown field such as
\texttt{persoan:} without any warning.

\paragraph{Bucket 4: workflow errors.} An arrow expression points at a
member that is not a valid workflow node, or attempts to subscript into a
non-list.

\begin{lstlisting}[caption={Workflow error: \texttt{tools} is not an agent.},label={lst:err-wf}]
tool search { invoke = "pkg.search" }
workflow pipeline {
  search -> writer
}
\end{lstlisting}

\noindent Orca reports \texttt{[unexpected-expr] workflow step must be an
agent}. LangGraph does not distinguish nodes by role, so the same
structural error is expressible only by convention and fails at runtime
with a type error inside the node function body.

\begin{imgplaceholder}[title={Screenshot: VS Code diagnostics on a broken \texttt{.oc}}]
LSP-driven diagnostics showing \texttt{[undefined-ref]},
\texttt{[type-mismatch]}, and \texttt{[missing-field]} highlights inline
within the editor as the file is typed.
\end{imgplaceholder}

Table~\ref{tab:error-coverage} summarizes which of the four buckets each
stack catches at author time versus first execution.

\begin{table}[t]
  \centering
  \caption{Stage at which each error class is caught (AT = author time
  via static analysis, RT = run time). cagent YAML is parsed at runtime
  and performs no cross-field validation, so all four classes surface as
  runtime failures.}
  \label{tab:error-coverage}
  \small
  \begin{tabular}{@{}lcccc@{}}
    \toprule
    \textbf{Error class} & \textbf{Orca} & \textbf{LangGraph} & \textbf{CrewAI} & \textbf{cagent} \\
    \midrule
    Reference errors  & AT & RT & RT & RT \\
    Type mismatches   & AT & RT & RT & RT \\
    Schema violations & AT & RT & RT & RT \\
    Workflow errors   & AT & RT & RT & RT \\
    \bottomrule
  \end{tabular}
\end{table}

The asymmetry in Table~\ref{tab:error-coverage} is the qualitative claim
that the SLOC numbers alone cannot capture: the two-line difference
between Orca and cagent in Table~\ref{tab:loc-results} is purchased by a
four-class difference in static guarantees.

\subsection{Threats to Validity}
\label{sec:evaluation:validity}

Several factors constrain the strength of these results. The use-case
implementations are hand-written by the authors and may inadvertently
favour Orca; a team that routinely writes LangGraph might factor common
node patterns into helper functions and recover part of the boilerplate
deficit. SLOC further ignores the cognitive load of framework conventions,
the cost of reading third-party documentation, and the differential
maintenance burden of imperative code as requirements evolve. The
evaluation is also non-empirical in the user-study sense: no developers
other than the authors have been timed on equivalent tasks, so authoring
effort is inferred from structural measurements rather than observed.
Finally, the error-class coverage table describes categories of failures
and not their relative frequency in production agentic systems, which
remains an open question.

Despite these caveats, the combination of the LOC reduction (especially
its super-linear growth across UC1--UC4), the structural asymmetry
between Orca and cagent on error coverage, and the side-by-side
comparison in Figure~\ref{fig:uc2-sidebyside} supports the two central
claims of the paper: the declarative DSL meaningfully reduces authoring
effort on non-trivial agentic systems and shifts classes of
misconfiguration that are traditionally runtime failures into compile-time
diagnostics.

% ══════════════════════════════════════════════════════════════════════════════
\section{Conclusion and Future Work}
\label{sec:conclusion}
% ══════════════════════════════════════════════════════════════════════════════

This paper introduced Orca, a declarative domain-specific language for AI
agent orchestration, and validated it along three complementary axes. The
design (Section~\ref{sec:design}) promotes agent configurations, tool
bindings, workflow topologies, and structured schemas to first-class
language primitives, with a structural type system that unifies primitive
types and user-defined records under a single \texttt{BlockRef} kind. The
implementation (Section~\ref{sec:implementation}) realises the language as
a four-stage Go compiler with a self-bootstrapped schema file, a Pratt
expression parser that makes the \texttt{->} workflow operator and the
\texttt{|} union operator share a single grammar, error-tolerant recovery
that drives a Language Server Protocol integration, and a browser-based
\emph{Orca Studio} canvas that treats the textual source as its single
authoritative representation. The evaluation
(Section~\ref{sec:evaluation}) demonstrated on four use cases of
increasing complexity that Orca achieves mean source-line reductions of
roughly $27.5\%$ against Python/LangGraph and $44.0\%$ against CrewAI,
with the saving growing super-linearly as workflow complexity increases,
and that it statically catches all four classes of misconfiguration in
the error taxonomy that the imperative and YAML baselines surface only at
runtime.

Several extensions are worth pursuing in future work.
\begin{enumerate}[leftmargin=*]
  \item \textbf{Additional code-generation backends.} The code generator
        currently targets LangGraph. Emitting CrewAI Python or Docker
        cagent YAML from the same validated AST would decouple the DSL
        from any single runtime and position Orca as a portable
        interchange format for agentic systems.
  \item \textbf{Richer workflow primitives.} Conditional edges,
        loops, and parallel fan-out are expressible in LangGraph but not
        yet in Orca's \texttt{->} grammar. Extending the Pratt expression
        language with guarded-edge and iteration constructs would close
        this expressiveness gap while preserving static analysability.
  \item \textbf{Empirical user study.} The Section~\ref{sec:evaluation:validity}
        threats to validity can only be fully discharged by timing
        developers on equivalent UC1--UC4 tasks and comparing authoring
        time, defect rate, and subjective cognitive load across the four
        stacks, rather than inferring effort from SLOC.
  \item \textbf{Studio feature expansion.} Surfacing live compile
        diagnostics directly on the canvas, visual diffing of two
        \texttt{.oc} versions, and SVG export for documentation would
        further strengthen the canvas as a collaboration surface for
        non-engineer stakeholders.
  \item \textbf{Package-level imports.} First-class module boundaries
        would let schemas, tools, and reusable agent templates be shared
        across \texttt{.oc} files and across teams, which is a
        prerequisite for an ecosystem of reusable orchestration building
        blocks.
\end{enumerate}

Taken together, the results support the thesis that agent orchestration is
amenable to the same language-design toolkit that previous waves of
infrastructure-as-code applied to cloud deployment: moving configuration
from imperative scripts into a declarative, statically-validated surface
yields measurable authoring-effort and reliability dividends without
sacrificing expressiveness.

% ══════════════════════════════════════════════════════════════════════════════
% References
% ══════════════════════════════════════════════════════════════════════════════
\bibliographystyle{ACM-Reference-Format}
\bibliography{references}

% ══════════════════════════════════════════════════════════════════════════════
\appendix
\section{Diagnostic-Code Catalogue}
\label{sec:appendix:diagnostics}
% ══════════════════════════════════════════════════════════════════════════════

Table~\ref{tab:diagnostics} lists every diagnostic code produced by the
parser and analyzer stages of the current Orca compiler. Each code can be
individually silenced at source level with
\texttt{@suppress("\textit{code}")}. Codes are stable across releases and
form the contract between the compiler and editor integrations.

\begin{table}[!ht]
  \centering
  \small
  \caption{Diagnostic codes emitted by the Orca compiler, grouped by
  bucket from Section~\ref{sec:evaluation:taxonomy}.}
  \label{tab:diagnostics}
  \begin{tabular}{@{}p{0.26\linewidth}p{0.19\linewidth}p{0.45\linewidth}@{}}
    \toprule
    \textbf{Code} & \textbf{Bucket} & \textbf{Condition} \\
    \midrule
    \texttt{syntax}            & ---           & Any parse-level error; used by the parser for malformed tokens and unrecovered recovery points. \\
    \midrule
    \texttt{undefined-ref}     & Reference     & Identifier not present in the symbol table. \\
    \texttt{unknown-member}    & Reference     & Member access on a block type that does not export that name. \\
    \midrule
    \texttt{type-mismatch}     & Type          & Field value's inferred type is incompatible with its schema-declared type. \\
    \texttt{invalid-subscript} & Type          & Non-integer expression used as a subscript on a list-typed value. \\
    \texttt{invalid-value}     & Type          & Field value is not a member of an enumerated allowed set. \\
    \midrule
    \texttt{missing-field}     & Schema        & Required field absent from a block body. \\
    \texttt{unknown-field}     & Schema        & Field name present but not declared by the block's schema. \\
    \texttt{duplicate-field}   & Schema        & Same field assigned twice within one block. \\
    \texttt{duplicate-block}   & Schema        & Two top-level blocks sharing the same identifier. \\
    \midrule
    \texttt{unexpected-expr}   & Workflow      & Expression used in a position where its kind is disallowed (e.g.\ non-agent as a workflow step). \\
    \texttt{unknown-provider}  & Backend       & Model provider string not supported by the selected code-generation backend. \\
    \texttt{unsupported-lang}  & Backend       & Raw-string language tag not recognized by the backend. \\
    \bottomrule
  \end{tabular}
\end{table}

\end{document}
